\chapter{Integración de modelos de lenguaje}

\section{Papel del LLM en el sistema}

El trabajo incorpora modelos de lenguaje como parte central del flujo. El LLM no se limita a redactar la respuesta final: primero interpreta la consulta y construye una intencionalidad estructurada; después genera código Python a partir de dicha interpretación y del mensaje original; finalmente interpreta las métricas devueltas por la ejecución.

Esta organización permite aprovechar la flexibilidad del modelo sin ejecutar sus salidas de forma ciega. Entre generación y ejecución existe una capa de control de seguridad que revisa el script antes de lanzarlo.

\section{Cliente compatible con OpenAI}

La integración se ha preparado mediante un cliente compatible con APIs tipo OpenAI. Esta decisión permite alternar entre proveedores sin modificar el workflow principal. El sistema utiliza variables de entorno para seleccionar perfil, modelo, URL base y clave de API.

En esta memoria, cuando se comentan resultados ya ejecutados, se hace referencia al modelo \nolinkurl{openai/gpt-oss-20b}~\cite{openai2025gptossmodelcard} servido mediante vLLM en el entorno universitario. Este modelo abierto de la familia gpt-oss se usa por su capacidad de razonamiento, seguimiento de instrucciones y generación de código Python. El nombre \texttt{university} queda reservado como identificador técnico del perfil de configuración dentro del código, no como nombre académico del modelo.

\section{Configuraciones evaluadas}

Se han preparado tres configuraciones principales:

\begin{itemize}
    \item \textbf{Groq con \nolinkurl{llama-3.3-70b-versatile}}: proveedor externo compatible con OpenAI. El modelo corresponde a Llama 3.3 70B Instruct y se utiliza por su razonamiento general, seguimiento de instrucciones y baja latencia en las pruebas~\cite{groqllama33docs,meta2024llama33modelcard,meta2024llama3herd}.
    \item \textbf{Gemini con \nolinkurl{gemini-2.5-flash}}: proveedor secundario para estudiar portabilidad y calidad comparativa. El modelo Flash de Gemini 2.5 se utiliza por su rapidez, capacidad de razonamiento y equilibrio entre latencia y calidad~\cite{google2025gemini25flashdocs,google2025gemini25report}.
    \item \textbf{\nolinkurl{openai/gpt-oss-20b} sobre vLLM universitario}: modelo abierto de la familia gpt-oss ejecutado en el servidor compatible con OpenAI Chat Completions habilitado en el entorno universitario. Se usa como modelo principal por su capacidad de razonamiento, seguimiento de instrucciones y generación de código~\cite{openai2025gptossmodelcard}.
\end{itemize}

\section{Gestión de errores de API}

El diseño registra errores de red, cuota, indisponibilidad, credenciales o formato de respuesta. Si una llamada LLM falla, el sistema devuelve un error explícito para que la evaluación no mezcle resultados reales de modelo con respuestas programáticas.

\section{Gestión de credenciales}

Las claves de Groq con \nolinkurl{llama-3.3-70b-versatile}, Gemini con \nolinkurl{gemini-2.5-flash} y del servidor vLLM universitario que sirve \nolinkurl{openai/gpt-oss-20b} se gestionan mediante variables de entorno en el fichero local \texttt{.env}. Este archivo no se versiona y no debe compartirse. El repositorio incluye configuraciones de ejemplo, pero una ejecución real con API solo es posible cuando el entorno local contiene una clave válida.



